\section{Optimasi Klasik SLSQP (\textit{Sequential Least Squares Programming})}

Sebagai standar pembanding (\textit{benchmark}) terhadap performa algoritma kuantum, penelitian ini menggunakan metode optimasi klasik \textit{Sequential Least Squares Programming} (SLSQP) untuk menyelesaikan masalah Markowitz kontinu. SLSQP merupakan algoritma optimasi iteratif yang dirancang untuk menangani masalah pemrograman non-linear dengan batasan kesamaan (\textit{equality constraints}) maupun batasan pertidaksamaan (\textit{inequality constraints}) secara efisien. Algoritma ini bekerja dengan merumuskan serangkaian sub-masalah \textit{Quadratic Programming} (QP) pada setiap langkah iterasi menggunakan aproksimasi deret Taylor orde kedua terhadap fungsi Lagrangian.

Proses optimasi dimulai dengan inisialisasi bobot portofolio awal yang seragam (\textit{uniform weights}) serta aproksimasi matriks Hessian sebagai matriks identitas. Pada setiap iterasi, gradien dari fungsi objektif dihitung dan matriks Hessian diperbarui menggunakan skema pembaruan BFGS (\textit{Broyden-Fletcher-Goldfarb-Shanno}) berdasarkan perubahan posisi dan gradien dari langkah sebelumnya. Arah pencarian optimal ditentukan dengan menyelesaikan sub-masalah QP yang mematuhi linierisasi dari batasan portofolio, yakni total bobot harus setara dengan satu dan setiap bobot berada dalam rentang $[0, 1]$. Algoritma ini dipastikan berkonvergensi menuju titik minimum lokal yang, untuk fungsi Markowitz yang bersifat konveks, juga merepresentasikan solusi optimal global. Rincian operasional dari algoritma SLSQP yang diimplementasikan disajikan pada Algoritma \ref{alg:slsqp_optimizer}.

\begin{algorithm}[ht]
\caption{Optimasi Klasik Markowitz menggunakan SLSQP}
\label{alg:slsqp_optimizer}
\begin{algorithmic}[1]
\Procedure{SLSQPOptimizer}{$\mu, \Sigma, \gamma$}
    \State $w_0 \gets \text{InitialWeights}(1/N)$ \Comment{Inisialisasi Seragam}
    \State $B_0 \gets \mathbf{I}$ \Comment{Inisialisasi Hessian Identitas}
    \For{$k = 0, 1, 2, \dots$ until converged}
        \State $g_k \gets \nabla f(w_k) = \gamma \Sigma w_k - \mu$ \Comment{Hitung Gradien}
        \State Solve Sub-QP: $\min_{d} (g_k^T d + \frac{1}{2} d^T B_k d) \text{ s.t. } \mathbf{1}^T d = 0, \text{ bounds}$
        \State $\alpha_k \gets \text{LineSearch}(w_k, d)$ \Comment{Cari ukuran langkah}
        \State $w_{k+1} \gets w_k + \alpha_k d$ \Comment{Update Bobot}
        \State $s_k \gets w_{k+1} - w_k, y_k \gets g_{k+1} - g_k$
        \State $B_{k+1} \gets \text{UpdateBFGS}(B_k, s_k, y_k)$ \Comment{Update Aproksimasi Hessian}
    \EndFor
    \State \Return $w_{final}$
\EndProcedure
\end{algorithmic}
\end{algorithm}
